\question[7]
Lab JOP (es.\ 05). Il binario definisce i gadget \lstinline|g_bin|/\lstinline|g_sh|/\lstinline|g_exec| (funzioni \lstinline|naked| che terminano con un salto indiretto al \lstinline|dispatcher|), un \lstinline|dispatcher| e un trampolino \lstinline|jop_start|; stesso overflow di \lstinline|rop.c| (offset 76).
(a) Dove vive la dispatch table e come fa \lstinline|jop_start| a trovarla senza conoscere un indirizzo di stack? (b) Qual è il ``program counter'' della catena su i386/x86-64/arm64? (c) In che cosa differisce dal ROP e perché su AArch64 il JOP è più ``naturale''?
\begin{solution}
(a) La dispatch table vive \textbf{nel payload dell'overflow, subito dopo il return address} (un'area dati qualsiasi). L'overflow dirotta il \lstinline|ret| su \lstinline|jop_start|: al suo ingresso lo stack pointer punta \emph{già} al primo elemento della tabella, che viene copiato nel registro-PC della catena (\lstinline|mov %esp,%edi| / \lstinline|mov %rsp,%r12| / \lstinline|mov sp,x19|). Resta tutto \emph{relativo}: serve solo l'\textbf{offset}, nessun indirizzo assoluto (niente leak/corefile).\\
(b) Un registro dedicato, callee-saved e non di argomento (così non è sporcato quando i gadget chiamano \lstinline|system|): \textbf{EDI} su i386, \textbf{R12} su x86-64, \textbf{X19} su arm64. Il \lstinline|dispatcher| legge il prossimo gadget da quel registro, lo avanza e vi salta (\lstinline|jmp *%eax| / \lstinline|br x1|).\\
(c) Nel ROP il motore è \lstinline|ESP|+\lstinline|ret| con la catena \emph{sullo stack}; nel JOP è il \lstinline|dispatcher| + un registro, con la catena nella \emph{dispatch table}, e \textbf{nessun} \lstinline|ret| $\to$ elude i rilevatori dei \lstinline|ret|. Su AArch64 la catena usa \lstinline|br|/\lstinline|b| (salti a registro): non si presenta il self-loop delle funzioni foglia via \lstinline|ret| che complica il \emph{ROP} su arm64 (es.\ 04) $\to$ il JOP è il caso in cui AArch64 funziona meglio di x86 per il riuso di codice a più stadi.
\end{solution}
